\documentclass[book.tex]{subfiles}
\begin{document}

\chapter{Coarse Grained Potentials}
\label{chap:coarse}
\noindent\rule{11cm}{0.4pt} \\
\textbf{Prerequisites:}  \autoref{chap:molecular_hamiltonian},   \\
\textbf{Difficulty Level:} ** \\
\noindent\rule{11cm}{0.4pt}
\vspace{5mm}

Coarse grained potentials hold out the promise of being able to enable simulations at much larger scales. However, the usability of these techniques has been limited by the challenge of parameterizing coarse grained force fields. New differentiable programming techniques hold out the promise of making coarse graining much more effective in practice.

\section{Mathematical Framework}

Consider a system with $N$ atoms in a constant-volume NVT ensemble. Let the coordinates of each atom be given by

\begin{align}
q = (q_1,\dotsc, q_N) \in \mathbb{R}^{3N}
\end{align}
We define the coarse grained variables
\begin{align}
\xi(q) &= \xi_1(q),\dotsc,\xi_M(q)
\end{align}
to be a smaller set of parameters with $M \ll 3N$. Suppose that the original probability distribution of the system is given by
\begin{align}
p(q) &= \frac{1}{Z} e^{-\beta V(q)}
\end{align}
where $Z$ is the partition function and $V$ is the potential energy. We can then define the probability distribution for the coarse-grained variables by
\begin{align}
p(\xi) &= \frac{1}{Z} \int e^{-\beta V(q)} \delta(\xi(q) - \xi) dq
\end{align}
This equation effectively sums over all full resolution states that collapse to coarse grained state $\xi$. We can also define a potential energy and forces on the coarse grained states by
\begin{align}
U(\xi) &= -\frac{1}{\beta} \ln p(\xi) \\
F(\xi) &= - \nabla_{\xi} U(\xi)
\end{align}
Unlike the classical potential, we may not have a closed-form for the coarse grained potential $U$ but we can represent it as a differentiable program.

\section{Machine Learned Coarse Grained Potentials}

To define a differentiable coarse-grained potential, we need to define the coarse graining mechanism to extract $\xi$ and we need to define a parameterized representation for $U(\xi)$. Assuming we have a coarse-grained form selected, we can parameterize the potential by a network
\begin{align}
U_\theta(\xi)
\end{align}
If our coarse-graining scheme breaks the original system into coarse-grained "particles", we can write
\begin{align}
U_\theta(\xi) &= \sum_i U_{i,\theta}(\xi)
\end{align}
where we sum the per-particle contribution in the Behler-Parrinello style.

Training this model requires some subtlety since high resolution trajectories have to be transformed into coarse-grained settings before we can train with standard methods like force-matching.
\end{document}